\documentclass[11pt]{article}
\usepackage[a4paper,margin=1in]{geometry}\usepackage{amsmath,amssymb,amsthm,microtype}
\newtheorem{theorem}{Theorem}[section]\newtheorem{corollary}[theorem]{Corollary}
\title{A Variable-Rank Block-Power Criterion for Moving Quotient Tails}\author{Bin Wang}\date{July 2026}
\begin{document}\maketitle
\begin{abstract}
We replace the fixed-space convergence hypothesis of RH-269 by a direct
variable-rank statement.  The operators may act on changing quotient spaces,
the selected rank and block length may grow, and only trace-norm/operator-
norm power bounds are required.  A geometric block decomposition then gives
the full RH-246 logarithmic tail estimate beyond the moving head.  The
criterion is exact and conditional; no such uniform moving certificate is
currently archived.
\end{abstract}
\section{Hypotheses}
Let $C_\sigma$ act on a Hilbert space $H_\sigma$, let $m_\sigma\ge2$, and
assume uniformly
\[
 \|C_\sigma^{m_\sigma}\|_1\le K_\sigma,\quad
 \|C_\sigma^{m_\sigma}\|\le\eta_\sigma<1,\quad
 \|C_\sigma^r\|\le L_\sigma^r\ (0\le r<m_\sigma).
\]
\section{Theorem}
\begin{theorem}[moving block tail]
For $n=\ell m_\sigma+r$, $0\le r<m_\sigma$, $n\ge m_\sigma$,
\[
 |\operatorname{Tr}C_\sigma^n|
 \le K_\sigma\eta_\sigma^{\ell-1}L_\sigma^r.
\]
Consequently, if $R>0$ and $\eta_\sigma R^{m_\sigma}<1$,
\[
 \sum_{n\ge m_\sigma}\frac{|\operatorname{Tr}C_\sigma^n|R^n}{n}
 \le
 \frac{K_\sigma R^{m_\sigma}}{m_\sigma(1-\eta_\sigma R^{m_\sigma})}
 \sum_{r=0}^{m_\sigma-1}(L_\sigma R)^r.
\]
\end{theorem}
\begin{proof}
Factor $C^n=(C^m)^{\ell-1}C^mC^r$, use the ideal inequality for the trace,
and group the geometric series by residue classes.  The denominator
$n^{-1}\le(\ell m)^{-1}\le m^{-1}$ gives the displayed bound.
\end{proof}
\section{Route consequence}
If $m_\sigma\to\infty$, $\eta_\sigma R^{m_\sigma}\to0$, and
\[
 \frac{K_\sigma R^{m_\sigma}}{m_\sigma}
 \sum_{r=0}^{m_\sigma-1}(L_\sigma R)^r\longrightarrow0,
\]
then, with a uniform prefix budget, the moving tail vanishes.  A convenient
stronger condition is
\[
 \limsup_{\sigma\downarrow0}K_\sigma^{1/m_\sigma}R
 \max\{1,L_\sigma R\}<1.
\]
This is precisely the rank-growing replacement for a fixed limiting
quotient.  Supplying these bounds is an open operator problem; finite
twelfth-power contractions do not supply them.
\end{document}
